%! TEX program = xelatex
% Example: two-page extended abstract (Submission Abstract pattern).
% Illustrative content. Authors, numbers, and results are placeholders.
% The references are real publications and show the expected format.
\documentclass[a4paper]{article}
\usepackage[english]{babel}
\usepackage{parskip}
\usepackage{geometry}
\usepackage{amsmath,amssymb}
\usepackage{assets/d3-abstract}

\newcommand{\tittel}{Decision-Focused Learning for Battery Storage Scheduling}

% Footer: target venue and page limit
\renewcommand{\projektname}{\textbf{Extended abstract for OR 2026}\enspace|\enspace\textbf{Page limit: 2}}

% Page contract: two pages, so two template calls below
\renewcommand{\abstractpages}{2}

\begin{document}
\newgeometry{bottom=0.5cm, right=1.5cm, left=1.5cm, top=1.5cm}
\pagestyle{empty}

% ============================================================
% Page 1: self-contained core (claim, approach, headline result)
% ============================================================

% Submission abstracts carry authors and affiliation in the subtitle
\renewcommand{\undertittel}{A. Author, B. Author\enspace|\enspace Chair of Information Systems and Business Analytics, University of W\"urzburg}

\renewcommand{\innledning}{%
We train price forecasts for battery scheduling on the cost of the resulting
dispatch decisions. On 24 months of day-ahead market data, the approach lowers
normalized regret by \highlight{31\%} relative to a two-stage forecast with
identical model capacity.}

\renewcommand{\venstre}{%
\sectionhead{Motivation}

Operators of grid-scale batteries schedule charging and discharging one day
ahead, based on forecasts of hourly electricity prices. Forecast models are
commonly trained on mean squared error. The dispatch cost, however, depends on
the ranking of hours within a day, and errors in hours near the charging
threshold change the schedule most.

Decision-focused learning trains the forecast on the downstream objective.
Elmachtoub and Grigas [1] derive the convex SPO+ surrogate for linear
objectives, and Donti et al. [2] differentiate through a stochastic program.
Bertsimas and Kallus [3] study the related prescriptive setting with auxiliary
data. Mandi et al. [4] survey the field and provide benchmarks. We study the
storage scheduling problem, which adds inter-temporal state constraints.

\sectionhead{Problem and Approach}

Let $c \in \mathbb{R}^{24}$ denote hourly prices and $x$ the features
available at bidding time. The operator solves
\[
  z^*(c) = \arg\min_{z \in \mathcal{Z}} \; c^\top z,
\]
where $\mathcal{Z}$ encodes power limits and the state-of-charge dynamics.
A forecast $\hat{c} = f_\theta(x)$ incurs the regret
$c^\top z^*(\hat{c}) - c^\top z^*(c)$.
We minimize the SPO+ surrogate of this regret over $\theta$.

\sectionhead{Training Procedure}

\begin{enumerate}
\item Pretrain $f_\theta$ on mean squared error for ten epochs.
\item Solve the scheduling LP for each training day, warm-started from the
  previous epoch's basis.
\item Update $\theta$ with the SPO+ subgradient and repeat from step 2 until
  the validation regret stops improving.
\end{enumerate}

\sectionhead{Data}

Hourly day-ahead prices for the German bidding zone, January 2024 to December
2025, with load, wind, and solar forecasts published before gate closure.
}

\renewcommand{\hoyre}{%
\sectionhead{Headline Results}

\begin{tikzpicture}
\node[infobox] {\textbf{\color{d3-navy}\Large 24}\quad months of out-of-sample evaluation};
\end{tikzpicture}

\vspace{1ex}

\begin{tikzpicture}
\node[infobox] {\textbf{\color{d3-navy}\Large 58\%}\quad less solver time per epoch};
\end{tikzpicture}

\sectionhead{Comparison}

\begin{tabular*}{\linewidth}{@{\extracolsep{\fill}} l S[table-format=1.3] S[table-format=2.1] @{}}
\toprule
\textbf{Method} & {\textbf{Regret}} & {\textbf{RMSE}} \\
\midrule
% --- Comparison methods ---
Two-stage (MSE)      & 0.118 & 9.4 \\
Two-stage (quantile) & 0.109 & 9.9 \\
\midrule
% --- Ours ---
\textbf{DFL, linear} & 0.094 & 11.2 \\
\textbf{DFL, neural} & \highlight{0.081} & 10.7 \\
\bottomrule
\end{tabular*}

\vspace{1ex}

Regret is normalized by the optimal dispatch cost. RMSE in EUR/MWh. All models
use the same features and the same number of parameters per class.

\sectionhead{Contributions}

\begin{enumerate}
\item A decision-focused training procedure for storage scheduling with
  state-of-charge constraints.
\item A warm-start scheme that cuts the solver time per epoch by 58\%.
\item An evaluation on 24 months of day-ahead prices with a rolling
  out-of-sample protocol.
\end{enumerate}

\vspace{2ex}

\sectionhead{Evaluation Design}

Models are retrained monthly on a rolling 12-month window and evaluated on the
following month. Page 2 reports the results by season and battery duration.

\vspace{2ex}

\keypoint{Training on dispatch cost lowers regret by 31\% while the forecast
RMSE rises by 1.3 EUR/MWh.}
}

% d3-abstract-template.tex
% Page layout for D3 abstracts. One \input of this file produces one page.
% Usage: populate the slot variables, then % d3-abstract-template.tex
% Page layout for D3 abstracts. One \input of this file produces one page.
% Usage: populate the slot variables, then % d3-abstract-template.tex
% Page layout for D3 abstracts. One \input of this file produces one page.
% Usage: populate the slot variables, then \input{assets/d3-abstract-template}

% === HEADER ===
\setheader{\tittel}{\undertittel}

% === LEAD PARAGRAPH (skipped when \innledning is empty) ===
\ingressblock

% === FULL-WIDTH BLOCK (skipped when \helbredde is empty) ===
\fullwidthblock

% === BODY (two columns) ===
\begin{multicols}{2}
\raggedcolumns

{\small \venstre \par}

\columnbreak

{\small \hoyre \par}

\end{multicols}

% === FOOTER ===
\bunntekst

\resetslots

\pagebreak


% === HEADER ===
\setheader{\tittel}{\undertittel}

% === LEAD PARAGRAPH (skipped when \innledning is empty) ===
\ingressblock

% === FULL-WIDTH BLOCK (skipped when \helbredde is empty) ===
\fullwidthblock

% === BODY (two columns) ===
\begin{multicols}{2}
\raggedcolumns

{\small \venstre \par}

\columnbreak

{\small \hoyre \par}

\end{multicols}

% === FOOTER ===
\bunntekst

\resetslots

\pagebreak


% === HEADER ===
\setheader{\tittel}{\undertittel}

% === LEAD PARAGRAPH (skipped when \innledning is empty) ===
\ingressblock

% === FULL-WIDTH BLOCK (skipped when \helbredde is empty) ===
\fullwidthblock

% === BODY (two columns) ===
\begin{multicols}{2}
\raggedcolumns

{\small \venstre \par}

\columnbreak

{\small \hoyre \par}

\end{multicols}

% === FOOTER ===
\bunntekst

\resetslots

\pagebreak


% ============================================================
% Page 2: continuation page (no lead paragraph, full-width table)
% ============================================================

\renewcommand{\undertittel}{Experiments, Discussion, and References}

\renewcommand{\helbredde}{%
\sectionhead{Regret by Season and Battery Duration}

\begin{tabular*}{\linewidth}{@{\extracolsep{\fill}} l
  S[table-format=1.3] S[table-format=1.3] S[table-format=1.3] S[table-format=1.3]
  S[table-format=1.3] S[table-format=1.3] @{}}
\toprule
& \multicolumn{3}{c}{\textbf{2-hour battery}} & \multicolumn{3}{c}{\textbf{4-hour battery}} \\
\cmidrule(lr){2-4}\cmidrule(lr){5-7}
\textbf{Method} & {\textbf{Winter}} & {\textbf{Summer}} & {\textbf{All}} & {\textbf{Winter}} & {\textbf{Summer}} & {\textbf{All}} \\
\midrule
Two-stage (MSE)      & 0.131 & 0.102 & 0.118 & 0.097 & 0.078 & 0.088 \\
Two-stage (quantile) & 0.120 & 0.096 & 0.109 & 0.092 & 0.075 & 0.084 \\
\midrule
\textbf{DFL, linear} & 0.103 & 0.084 & 0.094 & 0.080 & 0.066 & 0.073 \\
\textbf{DFL, neural} & 0.088 & 0.073 & \highlight{0.081} & 0.069 & 0.058 & \highlight{0.064} \\
\bottomrule
\end{tabular*}
}

\renewcommand{\venstre}{%
\sectionhead{Experimental Protocol}

We use hourly day-ahead prices for the German bidding zone from January 2024
to December 2025, together with load, wind, and solar forecasts published
before gate closure. Models are retrained monthly on a rolling 12-month window
and evaluated on the following month. Each configuration runs with five random
seeds, and the table reports the mean.

The battery has a round-trip efficiency of 90\% and one full cycle per day.
We consider durations of two and four hours at a fixed power rating.

\sectionhead{Discussion}

The regret reduction is largest in winter, when the daily price spread is wide
and the ranking of evening hours determines most of the revenue. With a 4-hour
battery the schedule covers more hours per day, single ranking errors carry
less weight, and the gap between the methods narrows.

The decision-focused models have a higher RMSE than the two-stage models. They
shift accuracy toward the hours that determine the schedule, which is
consistent with the analysis of the SPO+ loss in [1].

\sectionhead{Computation}

\begin{tabular*}{\linewidth}{@{\extracolsep{\fill}} l S[table-format=3.0] S[table-format=1.1] @{}}
\toprule
\textbf{Setup} & {\textbf{Min/epoch}} & {\textbf{Factor}} \\
\midrule
Two-stage (MSE)     & 4   & 1.0 \\
DFL, cold start     & 31  & 7.8 \\
\textbf{DFL, warm start} & 13 & 3.2 \\
\bottomrule
\end{tabular*}

\vspace{1ex}

Training time per epoch on one GPU node. The factor is relative to the
two-stage pipeline.
}

\renewcommand{\hoyre}{%
\sectionhead{Scope}

The scheduling model is a linear program with perfect foresight of the
battery state. Degradation costs and intraday re-optimization are outside the
scope of this study. All results refer to a price-taking operator.

\sectionhead{Conclusion and Outlook}

Training price forecasts on dispatch cost improves battery scheduling
decisions on real market data at a moderate computational cost. Next, we add a
degradation term to the objective, which turns the scheduling problem into a
quadratic program, and we extend the evaluation to the intraday market.

\sectionhead{References}

\begin{reflist}
\item Elmachtoub, A. N., Grigas, P. (2022). Smart ``Predict, then Optimize''.
  \textit{Management Science} 68(1), 9--26.
\item Donti, P. L., Amos, B., Kolter, J. Z. (2017). Task-based End-to-end Model
  Learning in Stochastic Optimization. \textit{Advances in Neural Information
  Processing Systems} 30.
\item Bertsimas, D., Kallus, N. (2020). From Predictive to Prescriptive
  Analytics. \textit{Management Science} 66(3), 1025--1044.
\item Mandi, J., Kotary, J., Berden, S., Mulamba, M., Bucarey, V., Guns, T.,
  Fioretto, F. (2024). Decision-Focused Learning: Foundations, State of the
  Art, Benchmark and Future Opportunities. \textit{Journal of Artificial
  Intelligence Research} 80.
\end{reflist}
}

% d3-abstract-template.tex
% Page layout for D3 abstracts. One \input of this file produces one page.
% Usage: populate the slot variables, then % d3-abstract-template.tex
% Page layout for D3 abstracts. One \input of this file produces one page.
% Usage: populate the slot variables, then % d3-abstract-template.tex
% Page layout for D3 abstracts. One \input of this file produces one page.
% Usage: populate the slot variables, then \input{assets/d3-abstract-template}

% === HEADER ===
\setheader{\tittel}{\undertittel}

% === LEAD PARAGRAPH (skipped when \innledning is empty) ===
\ingressblock

% === FULL-WIDTH BLOCK (skipped when \helbredde is empty) ===
\fullwidthblock

% === BODY (two columns) ===
\begin{multicols}{2}
\raggedcolumns

{\small \venstre \par}

\columnbreak

{\small \hoyre \par}

\end{multicols}

% === FOOTER ===
\bunntekst

\resetslots

\pagebreak


% === HEADER ===
\setheader{\tittel}{\undertittel}

% === LEAD PARAGRAPH (skipped when \innledning is empty) ===
\ingressblock

% === FULL-WIDTH BLOCK (skipped when \helbredde is empty) ===
\fullwidthblock

% === BODY (two columns) ===
\begin{multicols}{2}
\raggedcolumns

{\small \venstre \par}

\columnbreak

{\small \hoyre \par}

\end{multicols}

% === FOOTER ===
\bunntekst

\resetslots

\pagebreak


% === HEADER ===
\setheader{\tittel}{\undertittel}

% === LEAD PARAGRAPH (skipped when \innledning is empty) ===
\ingressblock

% === FULL-WIDTH BLOCK (skipped when \helbredde is empty) ===
\fullwidthblock

% === BODY (two columns) ===
\begin{multicols}{2}
\raggedcolumns

{\small \venstre \par}

\columnbreak

{\small \hoyre \par}

\end{multicols}

% === FOOTER ===
\bunntekst

\resetslots

\pagebreak


\end{document}
